\section*{अध्याय \ref{ch:g12:satellites-kepler} --- उपग्रह और ग्रहों की गति}

\begin{solution}{exo:g12:satellites-kepler:1}
मंच पर: $F = GMm/R^2 = \num{3.98e17}/\num{4.06e13} \approx
\qty{9.8e3}{N}$। $r = \qty{6.77e6}{m}$ पर: $F \approx \qty{8.7e3}{N}$ --- यानी ज़मीन वाले मान का अब भी $89\%$। उसमें
बैठे लोग इसलिए तैरते हैं कि वे स्थानक के साथ-साथ \emph{गिर} रहे हैं, इसलिए
नहीं कि गुरुत्व चला गया है।
\end{solution}

\begin{solution}{exo:g12:satellites-kepler:2}
$v = \sqrt{GM/R} = \sqrt{\num{3.98e14}/\num{6.37e6}} \approx
\qty{7.9e3}{m/s}$; $T = 2\pi R/v \approx \qty{5.1e3}{s} \approx
\qty{84}{min}$। ठीक तोप के गोले वाली \qty{7.9}{km/s}: सतह छूती कक्षा तो न्यूटन
का वही शॉट \emph{है}।
\end{solution}

\begin{solution}{exo:g12:satellites-kepler:3}
$v \propto 1/\sqrt{r}$: $\sqrt 2 \approx 1.41$ से बँटा। $T \propto r^{3/2}$: $2^{3/2} \approx 2.8$ से गुणा। इनमें से किसी में $m$ नहीं
है: इसलिए नहीं।
\end{solution}

\begin{solution}{exo:g12:satellites-kepler:4}
$v = \sqrt{\num{3.98e14}/\num{6.77e6}} \approx \qty{7.67e3}{m/s}$; $T = 2\pi r/v \approx \qty{5.55e3}{s} \approx \qty{92}{min}$; प्रतिदिन $\num{86400}/5546 \approx 15.6$ कक्षाएँ।
\end{solution}

\begin{solution}{exo:g12:satellites-kepler:5}
सूर्य को एक नाभि पर रखते \omterm{def:g12:satellites-kepler:ellipse}{दीर्घवृत्त}; बराबर समय में बराबर क्षेत्रफल;
और सभी ग्रहों के लिए एक ही $T^2/a^3$। सबसे तेज़ \omterm{def:g12:satellites-kepler:ellipse}{उपसौर} पर, दूसरे नियम से
(छोटे बुहारे गए खंड को तेज़ बुहारना पड़ता है)। नहीं: \omterm{def:g10:universal-gravitation:satellite}{उपग्रह} का
द्रव्यमान कट जाता है --- हिसाब में केवल केंद्रीय द्रव्यमान आता है।
\end{solution}

\begin{solution}{exo:g12:satellites-kepler:6}
$T = \qty{2.36e6}{s}$: $M = 4\pi^2 r^3/(GT^2) =
\num{2.24e27}/371 \approx \qty{6.0e24}{kg}$ --- यानी \qty{5.97e24}{kg} के $1\%$ के भीतर।
\end{solution}

\begin{solution}{exo:g12:satellites-kepler:7}
गुरुत्व पृथ्वी के केंद्र की ओर ताका है, और वृत्तीय गति में \omterm{def:g11:forces-and-motion:netforce}{परिणामी बल}
वृत्त के केंद्र की ओर: इसलिए दोनों केंद्र एक ही हैं। पेरिस के अक्षांश-वृत्त के
ऊपर टँगा वृत्त धुरी पर केंद्रित होता, पृथ्वी के केंद्र पर नहीं --- यानी असंभव।
टँगे रहना केवल वहीं काम करता है जहाँ अक्षांश-वृत्त केंद्र पर केंद्रित
महावृत्त हो: यानी \emph{भूमध्य रेखा} के ऊपर।
\end{solution}

\begin{solution}{exo:g12:satellites-kepler:8}
$T = \qty{2.754e4}{s}$: $M = 4\pi^2 r^3/(GT^2) =
\num{3.26e22}/\num{5.06e-2} \approx \qty{6.4e23}{kg}$ --- यानी आँकड़ा-पत्र वाला मंगल।
\end{solution}

\begin{solution}{exo:g12:satellites-kepler:9}
पृथ्वी: $T^2/a^3 = \num{9.96e14}/\num{3.35e33} =
\qty{2.98e-19}{s^2/m^3}$। मंगल: $T = \qty{5.94e7}{s}$, $\num{3.52e15}/\num{1.18e34} = \qty{2.98e-19}{s^2/m^3}$ --- बराबर, जैसा तीसरा नियम माँगता है। फिर
$M = 4\pi^2/(G \times \num{2.98e-19})
\approx \qty{1.99e30}{kg}$।
\end{solution}

\begin{solution}{exo:g12:satellites-kepler:10}
$v = 2\pi r/T = \qty{1.02e3}{m/s}$; $a = v^2/r =
\qty{2.72e-3}{m/s^2}$; $g/60^2 = 9.81/3600 = \qty{2.73e-3}{m/s^2}$। चंद्रमा $60$ पृथ्वी-त्रिज्या दूर बैठा है, और उसका
पतन सेब वाले से $60^2$ गुना कमज़ोर है: यानी ठीक व्युत्क्रम वर्ग --- दोनों के
लिए एक ही नियम।
\end{solution}

\begin{solution}{exo:g12:satellites-kepler:11}
सीरीस: $T = 2.77^{3/2} = \sqrt{21.3} \approx \qty{4.6}{yr}$। हैली: $T = 17.8^{3/2} \approx \qty{75}{yr}$, इसलिए अगला \omterm{def:g12:satellites-kepler:ellipse}{उपसौर} लगभग $1986 + 75 \approx 2061$ में।
\end{solution}

\begin{solution}{exo:g12:satellites-kepler:12}
\qty{400}{km} पर: $v \approx \qty{7.67e3}{m/s}$; \qty{350}{km} ($r = \qty{6.72e6}{m}$) पर: $v \approx \qty{7.70e3}{m/s}$। \omterm{def:g10:inertia:inventory}{घर्षण} सचमुच
\omterm{def:g11:mechanical-energy:em}{यांत्रिक ऊर्जा} खींच लेता है --- कक्षा सिकुड़ जाती है; पर नीचे
उतरते समय गुरुत्व का कार्य कर्षण से हुई चाल की हानि से कहीं अधिक चुका
देता है: $E_m$ गिरता है जबकि $E_k$ चढ़ता है।
\end{solution}

\begin{solution}{exo:g12:satellites-kepler:13}
$T = \qty{3.024e5}{s}$: यानी $M = 4\pi^2 r^3/(GT^2) =
\num{1.65e31}/6.10 \approx \qty{2.7e30}{kg} \approx 1.4$ सूर्य। \emph{ग्रह} का द्रव्यमान पहुँच से बाहर है: वह
कक्षा के समीकरणों में से कट गया था।
\end{solution}

\begin{solution}{exo:g12:satellites-kepler:14}
$T = \qty{8.862e4}{s}$, $GM = \qty{4.28e13}{m^3/s^2}$: $r^3 = GMT^2/(4\pi^2) = \num{8.52e21}$, इसलिए $r \approx
\qty{2.04e7}{m}$ और मंगल की ज़मीन से $h = r - R \approx \qty{1.7e7}{m} \approx
\qty{17000}{km}$ ऊपर।
\end{solution}

\begin{solution}{exo:g12:satellites-kepler:15}
$v_a = r_p v_p / r_a = \num{8.77e10} \times \num{5.45e4} /
\num{5.25e12} \approx \qty{9.1e2}{m/s}$। चूँकि $r_a/r_p \approx 60$, धूमकेतु $60$ गुना दूर और $60$ गुना धीमा है: वह दशकों
तक अँधेरे में सुस्ताता है और कुछ हफ़्तों में सूर्य के पास से सरपट निकल जाता है।
\end{solution}

\begin{solution}{pb:g12:satellites-kepler:1}
\textbf{1.} $T = \qty{86164}{s}$: ज़मीन पृथ्वी के साथ \emph{तारों के सापेक्ष} हर
\omterm{def:g12:satellites-kepler:geo}{नाक्षत्र दिवस} में एक चक्कर घूमती है; और \qty{24}{h} वाला दिन उसमें वह एक
डिग्री प्रतिदिन जोड़ देता है जितना पृथ्वी सूर्य के चारों ओर आगे बढ़ जाती है।

\textbf{2.} वृत्तीय गति का \omterm{def:g11:forces-and-motion:netforce}{परिणामी बल} वृत्त के केंद्र की ओर ताका होता
है; और एकमात्र बल, गुरुत्व, पृथ्वी के केंद्र की ओर: इसलिए दोनों केंद्र
एक ही बिंदु हैं।

\textbf{3.} पेरिस के ऊपर टँगा वृत्त धुरी पर, केंद्र से उत्तर की ओर केंद्रित
होता --- जिसे 2 ख़ारिज कर देता है। इसलिए कक्षा \emph{भूमध्यरेखीय तल}
में ही होनी चाहिए।

\textbf{4.} \omterm{thm:g12:newtons-laws:second}{न्यूटन के दूसरे नियम} का \omterm{def:g10:refraction:angles}{अभिलंब} घटक: $m v^2/r = GmM/r^2$, इसलिए $v = \sqrt{GM/r}$
($m$ कट जाता है)।

\textbf{5.} $T = 2\pi r/v = 2\pi\sqrt{r^3/(GM)}$; वर्ग करने पर $T^2/r^3 = 4\pi^2/(GM)$।

\textbf{6.} $r = \left(GMT^2/4\pi^2\right)^{1/3} =
(\num{7.49e22})^{1/3} \approx \qty{4.22e7}{m}$।

\textbf{7.} $h = r - R = \num{4.22e7} - \num{6.37e6} \approx
\qty{3.58e7}{m} \approx \qty{35800}{km}$।

\textbf{8.} $v = 2\pi \times \num{4.22e7}/\num{86164} \approx
\qty{3.07e3}{m/s}$।

\textbf{9.} $\sqrt{\num{3.98e14}/\num{4.22e7}} \approx
\qty{3.07e3}{m/s}$ --- संगत।

\textbf{10.} त्रिज्याएँ एक ही: द्रव्यमान 4 में ही कट गया था। भारी
\omterm{def:g10:universal-gravitation:satellite}{उपग्रह} केवल प्रक्षेपण का बिल बदलता है --- उसी कक्षा के लिए
अधिक ईंधन।

\textbf{11.} $v = \sqrt{\num{3.98e14}/\num{6.77e6}} \approx
\qty{7.67e3}{m/s}$।

\textbf{12.} $T = 2\pi r/v \approx \qty{5.55e3}{s} \approx
\qty{92.4}{min}$।

\textbf{13.} $\num{86400}/5546 \approx 15.6$ कक्षाएँ: यानी दिन में लगभग $16$ सूर्योदय।

\textbf{14.} $r_{\text{भूस्थिर}}/r_{\text{ISS}} = \num{4.22e7}/\num{6.77e6}
= 6.23$; $6.23^{3/2} \approx 15.5 = \num{86164}/5546$ --- केप्लर का तीसरा नियम, वहीं का वहीं
सत्यापित।

\textbf{15.} स्थानक तेज़ है (\qty{3.07}{km/s} के मुक़ाबले \qty{7.67}{km/s}), फिर भी
\omterm{def:g12:satellites-kepler:geo}{भूस्थिर} वाले पर प्रति \omterm{def:g10:orders-of-magnitude:unit}{किलोग्राम} अधिक ऊर्जा लगी:
स्थितिज ऊर्जा में हुई चढ़ाई चाल की हानि से भारी पड़ती है।

\textbf{16.} $M_J = 4\pi^2 r^3/(G T^2)$।

\textbf{17.} $T = \qty{1.529e5}{s}$: $M_J =
\num{2.97e27}/1.56 \approx \qty{1.90e27}{kg}$।

\textbf{18.} $\num{1.90e27}/\num{5.97e24} \approx 318$ पृथ्वियाँ --- और सूर्य का लगभग $1/1000$।

\textbf{19.} $T_J = \qty{3.576e4}{s}$: $r^3 = GM_J T_J^2/(4\pi^2) =
\num{4.11e24}$, $r \approx \qty{1.60e8}{m}$ --- और \qty{4.22e8}{m} पर बैठा आयो $2.6$ गुना
ऊँचा उड़ता है और बृहस्पति के आसमान में सरकता दिखता है।

\textbf{20.} प्रसारक को ऊँचाई \qty{35800}{km} वाले भूमध्यरेखीय वृत्त पर खड़ा कीजिए,
जहाँ वह \qty{3.07}{km/s} से सवार रहता है और तश्तरी के निशाने से कभी नहीं हटता; और उसी
नियम ने, आयो का महीना खाकर, बृहस्पति को $\qty{1.90e27}{kg}$ --- यानी $318$ पृथ्वियाँ ---
मुफ़्त में तौल दिया।
\end{solution}
